% ------------------------------------------------------------------
%  Capitolo 1 — Non per la scuola, ma per la vita
%  Parte I
%  Ricerca di riferimento: ricerca 01 §1, §2.5, §3.4, §16; 02 §2.5, §5; 03 §3.4, §6
%  Voci bibliografiche principali: ebbinghaus1885, murre2015, bahrick1984,
%    oecd2024piaac, karpicke2009
%  Epigrafe: ✔ verificata sul testo (vedi ricerca 03 e registro, ricerca 00)
%  Scheda del capitolo: ricerca/06_indice-ragionato.md, capitolo 1
%
%  STATO: prima stesura del capitolo di prova (22/09/2026), circa
%  3.200 parole. Dati verificati: Seneca, Ep. 106, 11–12 (testo latino);
%  S-05 (Ebbinghaus 68/38); D-12 (PIAAC); S-24 (Bahrick 1984, abstract);
%  Karpicke et al. 2009 (84% / 55% / 11% / 1%, S-31); D-20 (PIRLS/TIMSS,
%  formulazione qualitativa). Ricordo personale dell'autore inserito
%  il 23/09/2026 (esame di ecofisiologia acquatica): nessun \segnaposto
%  residuo nel capitolo.
% ------------------------------------------------------------------
\chapter{Non per la scuola, ma per la vita}
\label{cap:vita}

\epigrafe[Impariamo non per la vita, ma per la scuola.]{Non vitae sed scholae discimus.}{Seneca, Epistulae ad Lucilium 106, 12}

\iniziale{C}{i sono frasi} che si imparano senza averle mai studiate.
Una di queste è \enquote{non scholae sed vitae discimus}: non impariamo
per la scuola, ma per la vita. L'hai forse letta sulla copertina di un
diario, in un discorso di inizio anno, nella dedica di una tesi; in più
di una città europea sta scritta sopra l'ingresso di un liceo. Suona
come una promessa, e insieme come un rimprovero rivolto a chi studia
solo per il voto: ricordati che tutto questo ti servirà.

La frase è latina, e di solito la si attribuisce a Seneca. È giusto, ma
solo a metà. Seneca l'ha scritta davvero, in una delle lettere
all'amico Lucilio, verso il 65 dopo Cristo. Soltanto che l'ha scritta al
contrario: \emph{non vitae sed scholae discimus}, \enquote{impariamo
non per la vita, ma per la scuola}. Non era un ideale: era una
constatazione amara. La massima che da secoli dovrebbe ispirare lo
studio era, in origine, la diagnosi del suo fallimento.

Questo libro comincia da qui, perché quella diagnosi è ancora valida. E
comincia da una domanda che quasi ogni studente, prima o poi, si è
fatto, magari la sera dopo un esame andato bene: \emph{a che cosa serve
studiare, se poi si dimentica tutto?}

\section{La frase rovesciata}

Conviene leggere la lettera per intero, perché il contesto cambia il
senso.\footnote{Seneca, \emph{Epistulae morales ad Lucilium} 106,
11--12. Le traduzioni dei testi antichi, qui e nel resto del libro,
sono dell'autore, salvo diversa indicazione.} Lucilio aveva chiesto a
Seneca una questione sottile, di quelle che appassionavano le scuole
filosofiche del tempo: se il bene sia o no un corpo. Seneca risponde con
pazienza, argomento dopo argomento; poi si ferma, quasi imbarazzato, e
commenta: \emph{latrunculis ludimus}, \enquote{stiamo giocando a
dama}.\footnote{I \emph{latrunculi} erano un gioco da tavolo romano,
simile alla dama o agli scacchi.} E aggiunge: \emph{in supervacuis
subtilitas teritur: non faciunt bonos ista sed doctos} -- \enquote{la
finezza si consuma in cose superflue: queste cose non rendono buoni,
ma dotti}.

Seguono le righe che ci interessano. \emph{Apertior res est sapere,
immo simplicior: paucis satis est ad mentem bonam uti litteris},
\enquote{essere saggi è cosa più chiara, anzi più semplice: bastano
poche lettere per una mente buona}. Ma noi, continua Seneca, sprechiamo
lo studio come sprechiamo tutto il resto: \emph{quemadmodum omnium
rerum, sic litterarum quoque intemperantia laboramus: non vitae sed
scholae discimus}. \enquote{Come di tutte le cose, così anche delle
lettere soffriamo di intemperanza: impariamo non per la vita, ma per la
scuola.}

Due parole meritano attenzione. La prima è \emph{intemperantia}: non
la pigrizia, ma il suo contrario, l'eccesso disordinato. Seneca non
rimprovera a Lucilio di studiare poco, bensì di studiare molto e male,
accumulando sottigliezze che non cambiano nulla. La seconda è
\emph{schola}. Per un romano del I secolo la \emph{schola} non era la
scuola dell'obbligo, ma il luogo delle lezioni e delle dispute, dove i
maestri di filosofia e di retorica mettevano alla prova i loro allievi.
Studiare \enquote{per la scuola} voleva dire studiare per la gara, per
la figura, per la prossima discussione: per una prova, non per la vita.

Seneca non disprezzava lo studio: lo praticò fino alla fine e ne fece
il centro delle sue lettere. Disprezzava lo studio che si esaurisce
nella prova. Il suo lamento ha attraversato i secoli e, a un certo
punto, qualcuno lo ha capovolto. Non sappiamo con precisione chi sia
stato né quando; sappiamo che in età moderna la forma rovesciata è
diventata un motto delle scuole europee, e che in quella forma è
arrivata fino a noi.\footnote{La forma positiva non si trova in nessun
autore antico. Non è l'unico caso: \emph{docendo discimus} è una
formula moderna ricavata da un'altra lettera di Seneca (\emph{Ep.} 7, 8:
\emph{homines dum docent discunt}, \enquote{gli uomini, mentre
insegnano, imparano}), e \emph{repetita iuvant} non è una citazione
classica.} C'è
qualcosa di ironico in questo rovesciamento, e insieme qualcosa di
istruttivo: per trasformare una diagnosi in un augurio è bastato
invertire due parole. Per trasformarla in realtà, come vedremo, serve
qualcosa di più.

Se oggi Seneca entrasse in un'aula universitaria italiana, la sua
\emph{schola} avrebbe un nome preciso: l'appello. Pensa a come si
prepara, di solito, un esame. Si scopre la data, si calcola quante
pagine mancano, si divide il manuale per il numero di giorni. Nelle
ultime tre settimane si legge, si sottolinea, si riassume, si ripete;
negli ultimi tre giorni si ripassa fino a tardi. Poi l'esame: magari
va bene, magari benissimo. E dopo? Dopo il manuale torna sullo
scaffale, e non lo si riapre più. È uno studio fatto con serietà, a
volte con sacrificio. Ma è uno studio progettato per un giorno solo: il
giorno della prova. Studiare per l'appello è il nostro modo di
studiare \emph{scholae}.

\section{Che cosa resta dopo la scuola}

Fai una prova. Prendi il manuale di un esame che hai superato un anno
fa, o due. Le sottolineature sono tue, la grafia a margine è la tua; le
pagine che una volta sapevi quasi a memoria ti sembrano familiari,
eppure, se provi a spiegarne una a voce, ti accorgi che ne resta poco:
qualche parola, un'immagine, un'impressione. Come se quel libro lo
avesse studiato un'altra persona.

Ricordo ancora un mio esame di ecofisiologia acquatica, durante l'università. Studiai a fondo, produssi buoni riassunti al computer, memorizzai la tassonomia algale: presi 30 e lode. Il professore mi domandò perfino se la sua materia fosse di mio particolare interesse. Ero felicissimo. Ma dopo pochi mesi dimenticai tutto. Ad anni di distanza, l'unica nozione che ricordo di quell'esame è questa: esistono le \emph{diatomee}. Null'altro.

Forse anche tu hai avuto un'esperienza simile, dimenticare tutto o buona parte del materiale studiato. È uno dei fenomeni meglio documentati della
psicologia, e ha un inventore preciso. Nel 1885 un giovane studioso
tedesco, Hermann Ebbinghaus, pubblicò un libretto intitolato \emph{Sulla
memoria} \parencite{ebbinghaus1885}. Per anni aveva fatto da cavia a sé
stesso: imparava a memoria liste di sillabe senza senso -- scelte
apposta perché nessun significato lo aiutasse a ricordarle -- e poi
misurava quanto lavoro gli serviva per reimpararle dopo venti minuti,
un'ora, un giorno, un mese. Più lavoro serviva per reimparare, più
aveva dimenticato.

Il risultato è la celebre \emph{curva dell'oblio}
(figura~\ref{fig:oblio}). La dimenticanza non procede in modo regolare:
è rapidissima nelle prime ore e nei primi giorni, poi rallenta, e ciò
che sopravvive alle prime settimane tende a restare più a lungo. Nel
2015 due ricercatori olandesi, Jaap Murre e Joeri Dros, hanno ripetuto
l'esperimento di Ebbinghaus con lo stesso metodo, per circa settanta ore
di prove: la curva che hanno ottenuto coincide con quella del 1885, con
un dettaglio in più, una piccola risalita del ricordo dopo ventiquattro
ore, compatibile con il lavoro che il cervello svolge durante il sonno
\parencite{murre2015}. A più di un secolo di distanza, la curva
dell'oblio è ancora lì.

\begin{figure}[tbp]
  \centering
  \begin{tikzpicture}[x=0.95cm, y=2.6cm, font=\footnotesize]
    \draw[filetto!70, line width=0.3pt] (0,0) -- (6.4,0);
    \draw[filetto!70, line width=0.3pt] (0,0) -- (0,1.08);
    \foreach \x/\l in {0/0, 1/1, 2/2, 3/3, 4/4, 5/5, 6/6}
      \draw[filetto!70] (\x,0) -- (\x,-0.03) node[below] {\l};
    \foreach \y/\l in {0/0, 0.5/50, 1/100}
      \draw[filetto!70] (0,\y) -- (-0.08,\y) node[left] {\l};
    \node[below] at (3.2,-0.16) {\itshape settimane dallo studio};
    \node[rotate=90, above] at (-0.55,0.5) {\itshape \% ricordato};
    % curva senza ripresa
    \draw[inchiostro, line width=0.6pt, domain=0:6.2, samples=80]
      plot (\x, {0.22 + 0.78*exp(-1.3*\x)});
    % curva con riprese distanziate
    \draw[rubrica, line width=0.6pt, domain=0:1, samples=30]
      plot (\x, {0.22 + 0.78*exp(-1.3*\x)});
    \draw[rubrica, line width=0.6pt, dashed] (1,0.43) -- (1,1);
    \draw[rubrica, line width=0.6pt, domain=1:3, samples=40]
      plot (\x, {0.5 + 0.5*exp(-0.7*(\x-1))});
    \draw[rubrica, line width=0.6pt, dashed] (3,0.62) -- (3,1);
    \draw[rubrica, line width=0.6pt, domain=3:6.2, samples=40]
      plot (\x, {0.72 + 0.28*exp(-0.35*(\x-3))});
    \node[rubrica, anchor=west] at (3.9,0.95) {\itshape con due riprese};
    \node[anchor=west] at (3.6,0.14) {\itshape studiato una volta sola};
  \end{tikzpicture}
  \caption{La curva dell'oblio: andamento qualitativo, non dati di un
    esperimento. Ciò che si studia una volta sola si perde in fretta;
    ogni ripresa ben distanziata rende la curva più piatta (se ne
    parlerà nei capitoli~\ref{cap:recupero} e~\ref{cap:tempo}).}
  \label{fig:oblio}
\end{figure}

C'è però un secondo dato, meno noto e più incoraggiante. Negli anni
Ottanta lo psicologo americano Harry Bahrick volle sapere che cosa resta
di una lingua studiata a scuola dopo molti anni, anche se non la si usa
più. Mise alla prova 733 persone che avevano studiato spagnolo a scuola,
alcune pochi mesi prima, altre fino a cinquant'anni prima
\parencite{bahrick1984}. Il ricordo calava rapidamente nei primi tre-sei
anni; poi la curva si appiattiva e restava quasi immobile per decenni,
fino a trent'anni, prima di un ultimo calo. Bahrick
chiamò \emph{permastore}, \enquote{deposito permanente}, questa parte
della memoria che sopravvive quasi intatta. E osservò che quanto ne
entrava non dipendeva dal caso: dipendeva da quanto bene la lingua era
stata imparata all'inizio -- dal numero di corsi seguiti, dai voti
ottenuti.

Mettiamo insieme i due dati. Dimenticare è normale: la curva di
Ebbinghaus vale per tutti, per gli studenti brillanti come per quelli
distratti, per chi ha vent'anni come per chi ne ha sessanta. Ma non
tutto si perde allo stesso modo. Ciò che è stato imparato bene,
consolidato, ripreso nel tempo, può restare per tutta la vita. La
domanda giusta, allora, non è \enquote{perché dimentico?}, ma
\enquote{perché dimentico \emph{tutto}, e così in fretta?}. Alla prima
domanda risponde la biologia; alla seconda risponde il modo in cui
studiamo.

\begin{esercizio}
Prima di andare avanti, prendi un foglio e scegli un argomento che hai
studiato per un esame circa un anno fa. Senza aprire libri né appunti,
scrivi per cinque minuti tutto ciò che ne ricordi: definizioni, nomi,
esempi, collegamenti. Poi confronta con il manuale. Tieni il foglio: ti
servirà alla fine del capitolo, e di nuovo nel capitolo~\ref{cap:recupero}.
\end{esercizio}

Quello che vale per il singolo studente sembra valere, in grande, anche
per un Paese. Le indagini internazionali trovano i bambini italiani
della scuola primaria sopra la media internazionale, sia nella lettura
sia in matematica e scienze. Gli adulti, invece, sono tra gli ultimi.
Nell'ultima indagine dell'OCSE sulle competenze degli adulti tra i 16 e
i 65 anni, pubblicata nel dicembre 2024, l'Italia ha ottenuto 245 punti
nella comprensione di testi, contro una media OCSE di 260; 244 punti
nelle competenze matematiche, contro 263; 231 nella soluzione di
problemi, contro 251 \parencite{oecd2024piaac}. Sotto l'Italia, nella
comprensione dei testi, ci sono soltanto sei Paesi.

Bisogna leggere questi numeri con prudenza. Le prove per i bambini e
quelle per gli adulti non sono direttamente confrontabili; sui risultati
degli adulti pesano molte cose che con la scuola hanno poco a che fare:
il tipo di lavoro, l'età media della popolazione, la scarsissima
abitudine degli italiani a continuare a formarsi dopo gli studi. Non è
lecito concludere che la scuola italiana \enquote{non funziona}. Resta
però un fatto: nessuno dei dati disponibili descrive un sistema in cui
ciò che si impara si accumula e si consolida negli anni. Descrivono
piuttosto qualcosa di molto simile alla curva di Ebbinghaus applicata a
un intero Paese: molto studio, e poco che resta.

\section{Tre errori che si possono correggere}

Perché dimentichiamo quasi tutto ciò che studiamo? La risposta più
comune è anche la più scoraggiante: perché non siamo portati, perché
non abbiamo memoria, perché \enquote{non è la mia materia}. La ricerca
degli ultimi decenni ne dà un'altra, molto più utile. Le cause
principali sono tre, e nessuna ha a che fare con il talento. Sono
errori: e gli errori si possono correggere.

\subsection*{Il primo errore: fidarsi della sensazione}

Il primo è un errore di \emph{metacognizione}, cioè di giudizio su come
funziona la propria mente. Le strategie che \emph{sembrano} funzionare
non sono quelle che funzionano davvero. Rileggere un capitolo, ripassare
le pagine evidenziate, ricopiare un riassunto: sono attività piacevoli,
ordinate, rassicuranti. Alla seconda lettura tutto scorre, i concetti
sembrano chiari, le parole familiari; e quella familiarità viene
scambiata per conoscenza. Le strategie che funzionano davvero, invece
-- chiudere il libro e provare a ricordare, riprendere un argomento dopo
qualche giorno, mescolare esercizi diversi -- sono faticose, lente,
piene di vuoti e di errori, e danno la sgradevole impressione di non
stare imparando.

Non stupisce allora che gli studenti scelgano le prime. In
un'indagine condotta su studenti universitari americani, l'84\% ha
dichiarato di studiare rileggendo appunti o manuale, e il 55\% ha
indicato la rilettura come la propria strategia principale; solo
l'11\% si interrogava da solo per imparare, e appena l'1\% ne faceva la
strategia principale \parencite{karpicke2009}. In Italia dati di questo
tipo non esistono ancora, ma chiunque abbia frequentato un'aula studio
durante la sessione ha visto la stessa scena: evidenziatori, righelli,
pagine colorate. Perché la sensazione di sapere inganni in modo così
sistematico, e come difendersi, è l'argomento del
capitolo~\ref{cap:inganni}: forse il più importante del libro.

\subsection*{Il secondo errore: studiare per il giorno dopo}

Il secondo è un errore di \emph{progettazione}. Si studia per la prova
di domani, non per ricordare fra un anno: tutto il lavoro si concentra
nei giorni che precedono l'esame, e dopo l'esame l'argomento non viene
più ripreso. Anche questo lo aveva già misurato Ebbinghaus. Per
imparare una lista di sillabe in modo da ripeterla senza errori il
giorno dopo, gli servivano 68 ripetizioni fatte tutte in una volta;
lo stesso risultato si otteneva con appena 38 ripetizioni distribuite
nei tre giorni precedenti. La distribuzione nel tempo, concludeva, è
\enquote{decisamente più vantaggiosa} dell'accumulo in un'unica sessione
\parencite{ebbinghaus1885}. Quasi la metà del lavoro, per lo stesso
risultato. E il vantaggio, come vedremo nel capitolo~\ref{cap:tempo},
si fa sentire soprattutto quando si vuole ricordare a lungo.

Lo studio concentrato alla vigilia non è stupidità né pigrizia. È una
risposta razionale a un sistema che chiede di dimostrare una cosa in un
giorno preciso e poi non la chiede più. Se l'obiettivo è superare la
prova, studiare tutto negli ultimi giorni sembra perfino efficiente: il
materiale è fresco proprio quando serve. Il prezzo si paga dopo. Lo
paga lo studente, quando il corso del secondo anno presuppone quello del
primo e bisogna ristudiare da capo ciò che si era già studiato; lo paga
chi insegna, quando si accorge che i prerequisiti \enquote{fatti
l'anno scorso} non ci sono più. E lo paga il laureato, che esce
dall'università con un titolo pieno di esami e una memoria che ne
conserva una parte sorprendentemente piccola.

\subsection*{Il terzo errore: non sapere perché}

Il terzo è un \emph{vuoto di senso}. Se non sai perché studi una cosa,
la studi per il voto; e ciò che si studia per il voto, dopo il voto, si
dimentica, perché ha già esaurito il suo scopo. È esattamente il
lamento di Seneca: la \emph{schola} come fine, invece che come mezzo.
Non si tratta di trovare a ogni argomento un'utilità pratica immediata
-- molte delle cose più preziose che si imparano non servono a
\enquote{fare} qualcosa, ma a capire, a giudicare, a pensare meglio. Si
tratta di sapere che cosa lo studio cambierà in te, e che cosa ti
permetterà di fare per gli altri. Vedremo nel capitolo~\ref{cap:sapere}
che questa non è retorica: sapere perché si studia aiuta davvero a
studiare meglio, e ci sono esperimenti che lo mostrano.

\asterismo

Tre errori, dunque: uno di giudizio, uno di progetto, uno di senso. Non
sono colpe tue. Nessuno, probabilmente, ti ha mai insegnato come si
studia: non la scuola, che ha insegnato che cosa studiare, e quasi mai
l'università, che dà per scontato che tu lo sappia già. Lo studente
italiano impara a studiare da solo, per tentativi, copiando il metodo
del compagno di banco o quello del fratello maggiore. È un po' come se
si insegnasse a nuotare consegnando a ciascuno un costume e indicando
il mare. Qualcuno impara benissimo; molti restano a galla, con grande
fatica; alcuni rinunciano.

\section{Che cosa promette questo libro, e che cosa no}

Questo non è un libro di trucchi. Non troverai tecniche per leggere
mille parole al minuto, né il segreto per imparare un esame in sette
giorni, né un test per scoprire se sei una persona \enquote{visiva} o
\enquote{uditiva}. Molte promesse di questo genere sono semplicemente
false, e nel corso del libro vedremo perché.

Ti propone invece un metodo. Non un metodo inventato da noi: un
metodo noto, che la psicologia sperimentale ha messo alla prova per più
di un secolo, in centinaia di esperimenti, molti dei quali condotti
proprio su studenti universitari, e poi in aule vere, con programmi
veri ed esami veri. Si può riassumere in cinque verbi:
\emph{richiamare} dalla memoria invece di rileggere; \emph{distanziare}
lo studio nel tempo invece di concentrarlo; \emph{alternare} argomenti
e tipi di problemi; \emph{spiegare}, a sé stessi e agli altri;
\emph{collegare} ciò che si studia a ciò che si sa già. E un sesto, che
tiene insieme gli altri: sapere \emph{perché}.

Non è un metodo facile. Anzi, una delle scoperte più solide della
ricerca, e forse la più controintuitiva, è che una certa fatica è il
segno che si sta imparando: una fatica giusta, superabile, che rallenta
lo studio di oggi e costruisce il ricordo di domani. Chi cerca un modo
per studiare senza sforzo resterà deluso. Chi è disposto a faticare in
modo diverso -- spesso non di più, ma meglio -- troverà un modo di
studiare che fa risparmiare tempo, riduce l'ansia prima degli esami e,
soprattutto, lascia qualcosa dopo.

Il libro cercherà di essere onesto su ciò che si sa. La ricerca
sull'apprendimento non è tutta ugualmente solida: alcuni risultati sono
stati replicati centinaia di volte, in laboratorio e in classe; altri
sono promettenti ma ancora discussi; altri ancora, molto popolari, sono
semplicemente falsi. Ogni volta che sarà importante, lo dirò con
chiarezza, e citerò gli studi da cui provengono i dati: non per
sfoggiare erudizione, ma perché tu possa controllare. Un libro che
insegna a non fidarsi della sensazione di sapere non può chiederti di
fidarti sulla parola.

Ecco il percorso. La prima parte si chiede \emph{perché} si studia:
come si studia oggi all'università, e che cosa significa desiderare di
sapere. La seconda parte spiega \emph{come funziona la memoria}: che
cosa ne sapevano gli antichi, che cosa ne sa la scienza, e perché le
nostre intuizioni sull'apprendimento ci ingannano. La terza, il cuore
del libro, presenta \emph{il metodo}, un principio per capitolo, e lo
ricompone alla fine in una settimana di studio reale. La quarta parte
riguarda le \emph{condizioni} dell'apprendimento: il sonno, l'attenzione,
il telefono, l'intelligenza artificiale, la motivazione. La quinta porta
tutto questo nella vita \emph{universitaria}: il passaggio dalla scuola,
il semestre, gli esami scritti e orali, le diverse discipline, gli
studenti con DSA. La sesta, infine, parla a \emph{chi insegna}
all'università e a chi forma i futuri insegnanti: perché il metodo si
impara meglio quando anche le lezioni e gli esami sono pensati per farlo
funzionare.

Il libro, infine, applica a sé stesso ciò che insegna. Alla fine di ogni
capitolo troverai alcune domande: sono fatte per essere risposte a
libro chiuso, per iscritto o a voce, prima di controllare. Da questo
capitolo in poi, tra le domande ce ne sarà sempre almeno una su un
capitolo precedente, così che ciò che hai letto ritorni, a distanza di
tempo, quando ha cominciato a sbiadire. Il motivo lo capirai nel
capitolo~\ref{cap:recupero}. Per ora basta un consiglio: non saltarle.

\asterismo

Seneca chiedeva uno studio che rendesse migliori, non soltanto più
dotti. Duemila anni dopo, il suo lamento è diventato un motto che
abbiamo rovesciato, e poi dimenticato di mettere in pratica. Questo
libro prova a rovesciarlo una seconda volta, ma con i fatti: a mostrare
che imparare \emph{per la vita} non è un augurio da scrivere sopra la
porta di una scuola, bensì il risultato di un metodo preciso, che si
può imparare. Si comincia chiudendo questo libro e provando a ricordare.

\begin{daricordare}
Dimenticare è normale. Dimenticare \emph{tutto}, subito dopo l'esame, è
il segno di un metodo sbagliato, non di una testa sbagliata: e un metodo
si può cambiare.
\end{daricordare}

\begin{domande}
  \item Che cosa scriveva davvero Seneca, in quale contesto, e perché la
        sua frase è stata rovesciata?
  \item Quali sono i tre errori che fanno dimenticare ciò che si
        studia? Per ciascuno, prova a dire da dove nasce.
  \item Che cosa mostrano, insieme, la curva di Ebbinghaus e le
        ricerche di Bahrick sullo spagnolo imparato a scuola?
  \item Riprendi il foglio dell'esercizio di questo capitolo. Che cosa
        ricordavi dell'esame di un anno fa, e perché, secondo te,
        proprio quello?
\end{domande}

\begin{sintesi}
  \item Seneca scrisse \emph{non vitae sed scholae discimus} come un
        rimprovero: si studia per la prova, non per la vita. Il motto
        scolastico ne è il rovesciamento.
  \item Si dimentica in fretta ciò che si studia una volta sola (la
        curva dell'oblio); ma ciò che è stato imparato bene può restare
        per decenni.
  \item Le cause principali dell'oblio dopo l'esame sono tre errori
        correggibili: fidarsi della sensazione di sapere, studiare per
        il giorno dopo, non sapere perché si studia.
  \item Per correggerli esiste un metodo noto e fondato su prove:
        richiamare, distanziare, alternare, spiegare, collegare, e
        sapere perché.
\end{sintesi}

\finecapitolo
